% =========================================================================
% Ángel Sarrión — Curriculum Vitae
% Design: LaTeX Technical Precision + Synthwave Palette (Neon Pink & Cyan)
% =========================================================================
\documentclass[10pt,a4paper]{article}

\usepackage[utf8]{inputenc}
\usepackage[T1]{fontenc}
\usepackage{lmodern}
\usepackage[margin=1.4cm, top=1.4cm, bottom=1.4cm]{geometry}
\usepackage{xcolor}
\usepackage{hyperref}
\usepackage{enumitem}
\usepackage{titlesec}
\usepackage{tikz}
\usepackage{fontawesome5}

% --- Synthwave Color Palette ---
\definecolor{SynthPink}{HTML}{FF2A85}
\definecolor{SynthBlue}{HTML}{008CB4}
\definecolor{DarkVoid}{HTML}{0A0B16}
\definecolor{SlateGray}{HTML}{475569}

% --- Hyperref Setup ---
\hypersetup{
    colorlinks=true,
    linkcolor=SynthBlue,
    urlcolor=SynthBlue,
    pdftitle={Ángel Sarrión --- Senior Business Analyst Lead y Especialista en TI},
    pdfauthor={Ángel Sarrión}
}

% --- Section Styling ---
\titleformat{\section}
  {\large\bfseries\color{SynthPink}}
  {}
  {0em}
  {}[{\color{SynthBlue}\titlerule[1.2pt]}]

\titlespacing*{\section}{0pt}{12pt}{8pt}

\begin{document}
\pagestyle{empty}

% --- Header with Vector Icons (Matching Screenshot) ---
\begin{center}
    {\Huge\textbf{\scshape Ángel Sarrión}} \\[4pt]
    {\large\color{SynthBlue}\textbf{Senior Business Analyst Lead y Especialista en TI}} \\[6pt]
    {\small
        \faEnvelope~\href{mailto:angelsarrion@gmail.com}{angelsarrion@gmail.com} \ $\cdot$ \ 
        \faMapMarker*~{Albacete / Madrid, España \textbullet{} Experiencia en Reino Unido (Milton Keynes)} \\[3pt]
        \faLinkedin~\href{https://www.linkedin.com/in/angelsarrion}{LinkedIn} \ $\cdot$ \ 
        \faGithub~\href{https://github.com/angelsarrion}{GitHub} \ $\cdot$ \ 
        \faGlobe~\href{https://angelsarrion.github.io/ASCV/}{Portfolio}
    }
\end{center}

\vspace{-4pt}

% --- Section: Summary ---
\section{Perfil Ejecutivo}
\noindent Líder tecnológico con más de 18 años diseñando y entregando soluciones de TI de misión crítica e ingeniería de procesos en banca corporativa. Especialista en motores de decisión de riesgo crediticio, admisión automatizada e integración de sistemas a gran escala en Santander UK y Europa. Trayectoria demostrada alineando la visión estratégica de negocio con una ejecución técnica robusta bajo normativas estrictas y máxima disponibilidad operativa.

% --- Section: Top Projects ---
\section{Iniciativas Estratégicas Destacadas}
\noindent\textbf{\large Risk Decision Service (RDS)} \hfill {\color{SynthBlue}\textbf{2021 -- 2025}} \\
{\color{SynthPink}\textbf{Líder Técnico y Especialista Senior de TI $\cdot$ Santander UK}} \\
\vspace{2pt}
Motor central de decisión de riesgo crediticio en tiempo real para las evaluaciones de financiación particular y comercial en Santander UK, bajo estricto cumplimiento regulatorio y SLAs de baja latencia.
\begin{itemize}[leftmargin=1.5em, itemsep=1.5pt, topsep=2pt]
    \item Lideró el rediseño arquitectónico integral de los sistemas de evaluación de riesgo crediticio, incrementando precisión de scoring, rendimiento y trazabilidad regulatoria.
    \item Coordinó releases ágiles continuas de BAU y resolución de incidencias, asegurando evolutivos constantes sin disrupción en la operativa bancaria.
    \item Soporte técnico de guardia (On-Call) 24/7 para incidencias de alta severidad, garantizando máxima estabilidad y resiliencia en producción.
    \item Dirección técnica de squads multidisciplinares en contratos de API, modelos de datos e interdependencias de sistemas.
\end{itemize}
\noindent\textit{\footnotesize \textbf{Tech Stack:} Credit Risk Engines, REST APIs, Agile / Scrum, Jira \& Confluence, DevOps, Change Authority, Dynatrace, Data Modeling}
\vspace{10pt}
\noindent\textbf{\large Proyecto Spectrum} \hfill {\color{SynthBlue}\textbf{2018}} \\
{\color{SynthPink}\textbf{Analista Líder de Negocio y Sistemas $\cdot$ Santander UK}} \\
\vspace{2pt}
Iniciativa estratégica insignia del programa Santander DIGITAL: incorporación digital automatizada de principio a fin y concesión de tarjetas de crédito para clientes de banca de empresas.
\begin{itemize}[leftmargin=1.5em, itemsep=1.5pt, topsep=2pt]
    \item Diseñó la arquitectura funcional y técnica de bajo nivel que unió los canales digitales de clientes con los motores de scoring de riesgo en back-office.
    \item Integró IBM API Gateway (DataPower/Connect) para una orquestación de APIs empresariales de alto rendimiento y máxima seguridad.
    \item Implantó flujos de ingesta continua hacia el Data Lake en tiempo casi real (NRT) para analítica de riesgo y reporting regulatorio.
    \item Desplegó pipelines de CI/CD en entornos PaaS de nueva generación, fijando el estándar de rapidez de entrega y calidad en el banco.
\end{itemize}
\noindent\textit{\footnotesize \textbf{Tech Stack:} IBM API Gateway, Near Real-Time Data Lake, CI/CD, PaaS, DevOps, Business Banking, Onboarding Automation}
\vspace{10pt}
\noindent\textbf{\large Proyecto Risk PCSM} \hfill {\color{SynthBlue}\textbf{2016 -- 2017}} \\
{\color{SynthPink}\textbf{Analista Líder de Sistemas $\cdot$ Santander UK}} \\
\vspace{2pt}
Rediseño y modernización completa de toda la arquitectura de Riesgo de Crédito Minorista en Santander UK para cumplir con nuevas exigencias del mercado y automatizar la suscripción.
\begin{itemize}[leftmargin=1.5em, itemsep=1.5pt, topsep=2pt]
    \item Dirigió el análisis técnico y la entrega de los flujos de evaluación automática de riesgo para Tarjetas de Crédito, Cuentas Personales y Financiación a Empresas.
    \item Transformó los sistemas legados de admisión en arquitecturas modulares por servicios, reduciendo el tiempo de concesión de días a segundos.
    \item Alineó directores de Políticas de Riesgo, Arquitectos de Soluciones y desarrolladores mediante especificaciones vivas en Confluence y Jira.
    \item Implantó estrategias de pruebas de regresión automatizadas con DevOps, asegurando pases a producción sin incidencias.
\end{itemize}
\noindent\textit{\footnotesize \textbf{Tech Stack:} Retail Credit Risk, Admission Systems, Agile Delivery, Jira \& Confluence, Solution Architecture, Automated Underwriting}
\vspace{10pt}


% --- Section: Professional Experience ---
\section{Trayectoria Profesional}
\noindent\textbf{\large Santander Digital Services} \hfill {\color{SynthBlue}\textbf{Jun 2025 -- Actualidad}} \\
{\color{SynthPink}\textbf{Senior Business Analyst Lead}} $\cdot$ {\color{SlateGray}Boadilla del Monte, Madrid / Albacete, España}
\begin{itemize}[leftmargin=1.5em, itemsep=1.5pt, topsep=2pt]
    \item Define y refina historias de usuario en Jira con criterios de aceptación rigurosos, alineando las entregas de desarrollo con objetivos de negocio y cumplimiento.
    \item Coordina la actividad técnica de múltiples squads, resolviendo dudas de diseño y manteniendo una alta cadencia de entrega en todo el ciclo de vida.
    \item Estructura el alcance de pruebas y criterios de validación con los equipos de QA, acelerando la certificación de versiones automáticas y exploratorias.
    \item Colabora con stakeholders estratégicos para traducir requerimientos de negocio en soluciones técnicas escalables y viables para producción.
    \item Gestiona dependencias transversales entre plataformas y coordina con Change Authorities la mitigación de riesgos de integración y gobierno de cambios.
    \item Participa activamente en gobernanza de releases como DevOps, supervisando pases a producción y validación operativa.
\end{itemize}
\\\noindent\textit{\footnotesize \textbf{Stack:} Jira \& Confluence, Agile / Scrum, DevOps, Change Authority, Technical Architecture, Enterprise Integration, QA Automation}
\vspace{10pt}
\noindent\textbf{\large Santander UK} \hfill {\color{SynthBlue}\textbf{Ene 2021 -- May 2025}} \\
{\color{SynthPink}\textbf{Especialista Senior en TI (Líder Técnico RDS)}} $\cdot$ {\color{SlateGray}Milton Keynes, Reino Unido}
\begin{itemize}[leftmargin=1.5em, itemsep=1.5pt, topsep=2pt]
    \item Liderazgo técnico y rediseño arquitectónico del Risk Decision Service (RDS) que gobierna las evaluaciones de Riesgo de Crédito en tiempo real en Santander UK.
    \item Coordinó releases ágiles de BAU, triaje de defectos y optimización continua de rendimiento, asegurando operativa ininterrumpida en servicios de préstamo.
    \item Prestó soporte especializado de guardia (On-Call) 24/7, garantizando máxima disponibilidad en la infraestructura crítica de scoring.
    \item Dirigió equipos técnicos multidisciplinares, estableciendo estándares de integración y contratos de datos para los motores de decisión.
\end{itemize}
\\\noindent\textit{\footnotesize \textbf{Stack:} Risk Decision Service (RDS), Agile Releases, Dynatrace, Data Modeling, Incident Support (On-Call), Credit Risk Systems}
\vspace{10pt}
\noindent\textbf{\large Santander UK} \hfill {\color{SynthBlue}\textbf{Abr 2012 -- Dic 2020}} \\
{\color{SynthPink}\textbf{Analista de Sistemas / Analista de Negocio y Sistemas}} $\cdot$ {\color{SlateGray}Milton Keynes, Reino Unido}
\begin{itemize}[leftmargin=1.5em, itemsep=1.5pt, topsep=2pt]
    \item Miembro clave en la unidad Credit Risk IT, especializándose en sistemas de Admisión y promocionando a Lead Analyst para programas estratégicos del banco.
    \item Lideró iniciativas capitales del programa Santander DIGITAL, destacando Project Spectrum (2018) para onboarding digital de tarjetas a empresas.
    \item Dirigió el análisis de sistemas del Risk PCSM Project (2016-2017), modernizando el riesgo minorista en tarjetas, cuentas y financiación comercial.
    \item Colaboró estrechamente con Arquitectos de Soluciones en diseños de alto nivel y elaboró especificaciones técnicas de bajo nivel y modelado de datos.
    \item Elaboró documentación funcional y técnica exhaustiva en Confluence y gestionó requerimientos en Jira asegurando la transferencia de conocimiento.
\end{itemize}
\\\noindent\textit{\footnotesize \textbf{Stack:} Credit Risk IT, Project Spectrum, PCSM Project, IBM API Gateway, PaaS, Near Real-Time Data Lake, Solution Architecture, Confluence \& Jira}
\vspace{10pt}
\noindent\textbf{\large Isban UK} \hfill {\color{SynthBlue}\textbf{Abr 2011 -- Abr 2012}} \\
{\color{SynthPink}\textbf{Analista Funcional}} $\cdot$ {\color{SlateGray}Milton Keynes, Reino Unido}
\begin{itemize}[leftmargin=1.5em, itemsep=1.5pt, topsep=2pt]
    \item Análisis funcional e integración para el Proyecto Armadillo: plataforma de venta de Seguros de Vida de Santander en alianza con AVIVA UK.
    \item Responsable de implantación, pruebas y soporte sobre tecnologías heterogéneas: J2EE, .NET, Biztalk y sistemas centrales COBOL.
    \item Gestión de soporte operativo BAU y resolución de incidencias para asegurar la continuidad de negocio en sucursales y canales de venta.
\end{itemize}
\\\noindent\textit{\footnotesize \textbf{Stack:} J2EE, .NET, Biztalk, COBOL, Insurance Platforms, BAU Support}
\vspace{10pt}
\noindent\textbf{\large Vector ITC Group} \hfill {\color{SynthBlue}\textbf{Sep 2007 -- Feb 2011}} \\
{\color{SynthPink}\textbf{Desarrollador de Software}} $\cdot$ {\color{SlateGray}Albacete, España}
\begin{itemize}[leftmargin=1.5em, itemsep=1.5pt, topsep=2pt]
    \item Desarrollador Java EE responsable del análisis, diseño y programación de módulos para la banca minorista de Crèdit Andorrà y aplicaciones web corporativas de Santander.
    \item Desarrollo de servicios backend para el proyecto Red.es, gestionando el hosting y administración de dominios .es para el Gobierno de España.
\end{itemize}
\\\noindent\textit{\footnotesize \textbf{Stack:} Java EE, Retail Banking, Web Applications, Red.es Domain Systems, SQL}
\vspace{10pt}


% --- Section: Skills ---
\section{Matriz de Competencias}
\begin{itemize}[leftmargin=1.5em, itemsep=2pt, topsep=2pt]
\item \textbf{Motores de Riesgo de Crédito y Banca:} Risk Decision Service (RDS), Retail Credit Risk (PCSM), Admission \& Underwriting, Scoring Model Integration, Credit Cards \& Commercial Lending, Regulatory Compliance\item \textbf{Arquitectura e Integración:} Solution Architecture, IBM API Gateway (DataPower / Connect), RESTful API Design \& Contracts, Near Real-Time (NRT) Data Lake, PaaS \& Cloud Runtimes, Enterprise System Integration\item \textbf{Sistemas y Observabilidad:} Dynatrace APM, IT Infrastructure \& Systems Management, Incident Management \& On-Call (24/7), Data Modeling \& SQL, Java EE / .NET / Biztalk / COBOL\item \textbf{Metodologías y Liderazgo:} Technical Squad Leadership, Agile \& Scrum Delivery, Jira \& Confluence Governance, Change Authorities \& Risk Mitigation, DevOps \& CI/CD Pipelines, Executive Stakeholder Alignment
\end{itemize}

% --- Section: Education & Languages ---
\section{Formación Académica e Idiomas}
\noindent\textbf{IES Leonardo da Vinci} \hfill {\textbf{2002 -- 2005}} \\
Técnico Superior en Desarrollo de Aplicaciones Informáticas (DAI) -- \textit{Análisis, Arquitectura de Software y Desarrollo de Aplicaciones}\\

\vspace{6pt}
\noindent \textbf{Español:} Nativo / Bilingüe, \textbf{Inglés:} Competencia Profesional Completa

\end{document}
